% Section: Hysteresis × N10 5-Seed Panel Replication (P7, iter 88)
% SYNTH (closes the iter-87 mint-recommendation: extend zvf-triage
% hysteresis to the N10 5-seed panel to certify the iter-87 result is
% not seed-specific). Pairs with iter-87 §sec:p7-iter87-hysteresis.
\subsection{Hysteresis \texorpdfstring{$\times$}{x} N10 5-Seed Panel Replication}
\label{sec:p7-hysteresis-n10}

The iter-87 hysteresis (\S\ref{sec:p7-iter87-hysteresis}) result on the
N2 four-method corpus reduces flips to 33\% of raw at matched yield. The
brief-mint-recommendation explicitly asks whether this result is
\emph{seed-robust} on an independent panel. We replicate on the N10
5-seed GRPO panel with 15 steps per seed (aggregate step-level $\mathrm{zvf}$
only, no per-prompt $k$ distribution). The replication uses the same
hysteresis filter as iter-87 with $\tau \in \{0.40, 0.50, 0.60, 0.70\}$
and $K_\text{up} = K_\text{dn} \in \{2, 3\}$ (8 non-baseline configs).
Three proxy outcome metrics are computed on each 15-step trajectory:
\emph{fires} (escalated-step count), \emph{flips} (state-change count),
and \emph{yield\_proxy} ($= n^{-1} \sum_{t\,:\,s_t=1} \mathrm{zvf}_t$,
the average ZVF over escalated steps).

\subsubsection{Headlines}

\paragraph{H1 (flip-flop hazard is REAL on N10 too).}
At $\tau = 0.60$, raw zvf-triage fires on 6--10/15 steps per seed and flips
7--10 times per trajectory. The mean raw flip-rate is
$9.0/15 = 0.6$ flips/step. Hysteresis@$K{=}2$ reduces this to
$\mathbf{2.0/15 = 0.13}$ flips/step --- a $\mathbf{5.0\times}$ reduction
in flip-rate. Paired-seed bootstrap median flip-ratio $\mathbf{= 0.157}$,
95\% CI $[0.111,\, 0.212]$, excludes $1.0$.

\paragraph{H2 (Pareto dominance at $\tau = 0.50$, $K = (2,2)$ on N10).}

\begin{table}[h]
\centering
\small
\begin{tabular}{rccc}
\toprule
$\tau$ & flip-ratio (median) & 95\% CI & yield-retention (median) \\
\midrule
0.40 & 0.402 & $[0.174,\, 0.700]$ & \textbf{1.009} $[0.969,\, 1.049]$ \\
\textbf{0.50} & \textbf{0.400} & $\mathbf{[0.174,\, 0.740]}$ & \textbf{1.009} $[0.970,\, 1.049]$ \\
0.60 & 0.157 & $[0.111,\, 0.212]$ & 0.814 $[0.519,\, 1.014]$ \\
0.70 & 0.157 & $[0.000,\, 0.357]$ & 0.363 $[0.000,\, 0.796]$ \\
\bottomrule
\end{tabular}
\caption{Paired-seed bootstrap on (flip-ratio, yield-retention) for
$K=(2,2)$ across N10 5-seed GRPO panel. The $\tau \in \{0.40, 0.50\}$
cells are Pareto-dominant: flip-ratio CI excludes 1.0 and yield-retention
CI brackets 1.0. $\tau = 0.50$ is the observed operating point on this panel.}
\label{tab:p7-hyst-n10}
\end{table}

\paragraph{H3 (cross-panel replication: N2 and N10 agree).}

\begin{table}[h]
\centering
\small
\begin{tabular}{lcc}
\toprule
panel & $\tau=0.50, K=(2,2)$ flip-ratio [CI] & yield-retention [CI] \\
\midrule
N2 (4 methods, 40 steps, 16 prompts) & 0.31--0.33 [0.07, 0.65] & 1.13--1.66 [0.52, 2.38] \\
N10 (5 seeds, 15 steps)             & 0.40 [0.17, 0.74]      & 1.009 [0.97, 1.05] \\
\bottomrule
\end{tabular}
\caption{The qualitative Pareto-dominance of hysteresis@$\tau{=}0.50,
K{=}(2,2)$ is \textbf{replicated across panels and seeds}.}
\label{tab:p7-hyst-cross-panel}
\end{table}

The cross-panel agreement on the **direction** (hysteresis reduces flips
to $\le 40\%$ of raw, yield preserved at $\ge 95\%$ level) is the
operational result. The quantitative difference --- N2 yield-retention
above 1.0 --- is a \emph{gift-specific} effect (iter-87 H5): gift's steep
local ZVF drops create over-escalation in raw zvf-triage; hysteresis
filters these out and \emph{raises} gift's effective yield. N10 is
GRPO-only, so the gift effect cannot transfer.

\subsubsection{Prospective candidate policy}
\begin{itemize}
\item \textbf{Hysteresis@$\tau{=}0.50, K_\text{up}{=}K_\text{dn}{=}2$}
      is a pre-registered candidate policy, not a default filter; it requires
      matched-budget held-out evaluation on each target panel. The N10 panel
      (5 seeds × 15 steps) descriptively shows flip-ratio 0.40 of raw
      (CI $[0.17, 0.74]$ excludes 1.0), yield-retention 1.009 of raw
      (CI $[0.97, 1.05]$ brackets 1.0).
\item At $\tau = 0.60$ hysteresis is too aggressive on N10 (yield drops
      to 81\%); at $\tau = 0.70$ it is over-aggressive (yield 36\%).
      These descriptive comparisons make $\tau = 0.50$ a candidate to test,
      not an established operating point.
\end{itemize}

\subsubsection{Cross-paper coupling}
\begin{itemize}
\item \textbf{P7 iter-87 (\#103, N2 four-method)}: iter-88 is the
      \textbf{independent 5-seed replication} on N10. The
      qualitative Pareto-dominance survives panel and seed.
\item \textbf{P7 iter-79 (\#93, multi-trigger seed-robust)}: iter-79
      measured per-seed trigger-fire stability at the fire-count level;
      iter-88 measures the same at the hysteresis-filter level. Both
      extend the P7 family by per-seed statistical certification.
\item \textbf{P5P8-SYNTH JOB B (this iteration)}: per the brief, this
      iteration's SYNTH JOB drives the top brief-mint-recommendation
      to validated. Iter-87's mint listed 6 fresh veins; 4 of 6 were P7
      extensions; the +N10 5-seed test was the most prominent statistical
      generalizability concern.
\end{itemize}
